\chapter{Regler}

\subsection{Definition}
Ein Regler ist eine Vorrichtung, die eine Ausgangsgröße kontinuierlich misst und durch Regelung einer Eingangsgröße konstant hält oder in einen Sollwert nachführt.

\subsection{Regelkreis}
Der Regelkreis besteht aus folgenden Komponenten:
\begin{itemize}
    \item \textbf{Sollwert}: Gewünschter Zielwert
    \item \textbf{Istwert}: Aktuell gemessener Wert
    \item \textbf{Regler}: Vergleicht Soll- und Istwert
    \item \textbf{Stellglied}: Führt Regelung durch
    \item \textbf{Strecke}: Zu regelnder Prozess
\end{itemize}

\subsection{Reglervarianten}
Häufige Reglervarianten sind:
\begin{description}
    \item[P-Regler:] Proportionale Verstärkung
    \item[I-Regler:] Integrale Ausregelung
    \item[D-Regler:] Differentiale Vorausnahme
    \item[PID-Regler:] Kombination aller drei Anteile
\end{description}
